\title{DeepSeekMath: Pushing the Limits of Mathematical Reasoning in Open Language Models}

\begin{abstract}
Mathematical reasoning remains a formidable obstacle for language models, given its intricate and highly organized structure. In this study, we present DeepSeekMath~7B, which extends the pre-training of DeepSeek-Coder-Base-v1.5 7B utilizing 120B tokens related to mathematics, collected from Common Crawl, alongside data from natural language and programming code. DeepSeekMath~7B attains a remarkable 51.7\% accuracy on the MATH benchmark at the competition level, doing so without the use of external toolkits or ensembling strategies, and thus approaching the performance exhibited by models like Gemini-Ultra and GPT-4. When applying self-consistency across 64 model outputs, DeepSeekMath~7B attains a score of 60.9\% on the same benchmark. The exceptional mathematical reasoning skill of DeepSeekMath~is primarily the result of two elements: First, we effectively exploit the vast reserves of open-access web data via a rigorously developed data filtering pipeline. Second, we propose Group Relative Policy Optimization (GRPO), a novel modification of Proximal Policy Optimization (PPO), which not only strengthens the model’s mathematical reasoning proficiency but also improves PPO's memory efficiency.
\end{abstract}

\section{Introduction}

The advent of large language models (LLMs) has fundamentally changed how artificial intelligence systems address mathematical reasoning, enabling notable progress in quantitative reasoning benchmarks \citep{MATH} as well as geometry-focused evaluations \citep{trinh2024solving}. Additionally, such models are increasingly valuable for supporting human efforts to tackle challenging mathematical tasks \citep{tao}. Nevertheless, state-of-the-art systems such as GPT-4 \citep{gpt4} and Gemini-Ultra \citep{gemini} remain proprietary, and the currently available open-source alternatives continue to lag substantially in terms of performance.

This paper presents DeepSeekMath, a specialized language model for mathematics that demonstrates substantial improvements over existing open-source models and achieves results comparable to GPT-4 across academic evaluation sets. Central to this advancement is the construction of the DeepSeekMath Corpus, a massive and meticulously curated pre-training dataset encompassing 120B math-specific tokens. The data is sourced from Common Crawl (CC) and identified through a classifier utilizing fastText \citep{joulin2016fasttext}. Initially, the classifier is trained with positive samples from OpenWebMath \citep{openwebmath} and a broad array of negative samples from various web pages. This trained model is subsequently leveraged to uncover further positive examples from the CC, which are then subjected to human annotation for enhanced quality. The improved dataset is used to iteratively refine the classifier’s performance. Our evaluations demonstrate that the resulting corpus is of considerable quality, as evidenced by our foundational model, DeepSeekMath-Base 7B, achieving 64.2\% accuracy on GSM8K \citep{gsm8k} and 36.2\% on the MATH competition benchmark \citep{MATH}, surpassing the performance of Minerva 540B \citep{minerva}. Additionally, due to the corpus’s multilingual composition, the model shows notable gains on Chinese math benchmarks \citep{wei2023cmath,agieval}. We consider our methodology for processing mathematical datasets as a foundation for the wider research community, anticipating further advancements in this domain.

DeepSeekMath-Base is constructed by initializing from DeepSeek-Coder-Base-v1.5 7B \citep{deepseek-coder}, as we found that leveraging a model pre-trained on code yields superior results compared to starting with a standard large language model. Notably, we also observed that mathematical pre-training enhances performance not only in mathematics-specific domains but also on broader benchmarks such as MMLU \citep{mmlu} and BBH \citep{bbh}, suggesting a general uplift in reasoning capabilities in addition to mathematical proficiency.

Subsequent to pre-training, DeepSeekMath-Base undergoes instruction tuning on mathematical tasks, utilizing datasets for chain-of-thought \citep{cot}, program-of-thought \citep{pot,pal}, and tool-augmented reasoning \citep{tora}. This process yields DeepSeekMath-Instruct 7B, which outperforms all other 7B models and attains results on par with larger, instruction-tuned open-source models in the 70B parameter class.

In addition, we propose Group Relative Policy Optimization (GRPO), a reinforcement learning algorithm related to Proximal Policy Optimization (PPO) \citep{schulman2017proximal}. GRPO dispenses with the critic network and instead calculates the baseline from aggregated group performance, substantially reducing computational requirements. Employing GRPO exclusively with a subset of English instruction tuning data produces marked gains over DeepSeekMath-Instruct, including notable increases across both familiar domains (GSM8K: 82.9\% $\rightarrow$ 88.2\%, MATH: 46.8\% $\rightarrow$ 51.7\%) and transfer domains (e.g., CMATH: 84.6\% $\rightarrow$ 88.8\%) during the RL process. We further propose a cohesive framework that contextualizes Rejection Sampling Fine-Tuning (RFT) \citep{yuan2023scaling}, Direct Preference Optimization (DPO) \citep{dpo}, PPO, and GRPO as variations or simplifications of RL. Our framework facilitates comprehensive experimental analysis, including contrasts such as online versus offline learning, outcome versus process-based supervision, and comparisons between single-pass and iterative RL. Lastly, we elucidate how our RL methodology enhances the effectiveness of instruction-tuned models and outline prospective strategies for achieving even more robust RL through this unified conceptual lens.

\subsection{Contributions}
This work presents advancements in large-scale mathematical pre-training, as well as an in-depth investigation of reinforcement learning methodologies.

\textbf{Math Pre-Training at Scale}

\begin{itemize}[topsep=0pt]
    \item We demonstrate that the open-source Common Crawl dataset is a rich resource for mathematical information. Through the creation of a carefully curated data selection process, we assemble the DeepSeekMath Corpus—an extensive, high-quality dataset comprising 120B tokens from mathematically-relevant web content. This collection is nearly 7 times more expansive than the dataset used by Minerva \citep{minerva} and approximately 9 times larger than OpenWebMath \citep{openwebmath}.

    \item Our model, DeepSeekMath-Base 7B, when evaluated, achieves results on par with Minerva 540B \citep{minerva}. This demonstrates that sheer parameter count is not the predominant factor determining mathematical reasoning performance; well-designed smaller models, if trained on appropriately selected data, can reach comparable levels of capability.

    \item Insights gleaned from our math-oriented pre-training experiments indicate that preceding math pre-training with code-focused training leads to notable gains in mathematical problem-solving, whether or not auxiliary tools are employed. This finding provides partial resolution to the persistent query: \textit{does code training enhance reasoning skills?} Our results support the assertion, especially for mathematical domains.

    \item We observe that the widely adopted practice of leveraging arXiv papers during training, while popular among math models, fails to yield significant benefits across the diverse mathematical benchmarks assessed in our study.
\end{itemize}

\textbf{Exploration and Analysis of Reinforcement Learning}

\begin{itemize}[topsep=0pt]
    \item We present Group Relative Policy Optimization (GRPO), a novel and resource-efficient algorithm for reinforcement learning. Unlike Proximal Policy Optimization (PPO), GRPO eliminates the need for a critic model by instead utilizing group scores to estimate the baseline, leading to substantial savings in training computation.
    \item Our results show that applying GRPO to the instruction-tuning data of DeepSeekMath-Instruct yields notable improvements. In addition, the reinforcement learning phase also contributes to enhanced out-of-domain generalization.
    \item We propose a unified framework to analyze and relate existing techniques including RFT, DPO, PPO, and GRPO. Through a series of thorough experiments—ranging from online versus offline training, supervision of outcomes versus processes, to comparisons of single-turn versus iterative reinforcement learning—we systematically probe the foundational aspects of this overarching framework.
    \item Drawing on the insights from our unified perspective, we further investigate the drivers of reinforcement learning efficacy, and outline a set of promising directions for advancing large language model (LLM) reinforcement learning methods.
\end{itemize}

\subsection{Summary of Evaluations and Metrics}

\begin{itemize}[topsep=0pt]
    \item \textbf{English and Chinese Mathematical Reasoning}: We perform systematic evaluations of our models on a diverse set of mathematical reasoning datasets in both English and Chinese, spanning problems from elementary to collegiate levels. The English datasets used include GSM8K \citep{gsm8k}, MATH \citep{MATH}, SAT \citep{llemma}, OCW Courses \citep{minerva}, and MMLU-STEM \citep{mmlu}. For Chinese, evaluations are based on MGSM-zh \citep{mgsm}, CMATH \citep{wei2023cmath}, Gaokao-MathCloze \citep{agieval}, and Gaokao-MathQA \citep{agieval}. Our assessment covers both models' abilities to construct fully self-contained textual solutions (without relying on external tools) and to solve tasks via Python code.

    In the English evaluation, DeepSeekMath-Base demonstrates a performance level on par with the proprietary Minerva 540B model \citep{minerva}, and consistently outperforms open-source baselines such as Mistral 7B \citep{mistral} and Llemma-34B \citep{llemma}, regardless of their pre-training history, often by a substantial margin. For Chinese datasets, DeepSeekMath-Base delivers leading results, plausibly owing to the inclusion of high-quality multilingual pre-training data rather than the English-centric approach in previous studies \citep{minerva,llemma}. Through mathematical instruction tuning and reinforcement learning, DeepSeekMath-Instruct and DeepSeekMath-RL achieve notable advancements, including surpassing the 50\% accuracy threshold on the highly challenging competition-level MATH dataset—marking a first within the open-source community.

\item \textbf{Formal Mathematics}: The performance of DeepSeekMath-Base is assessed using the informal-to-formal theorem proving benchmark introduced by \citep{dsp_proof}, with the evaluation conducted on miniF2F \citep{minif2f} and Isabelle \citep{isabelle} selected as the proof assistant. Results indicate that DeepSeekMath-Base excels in few-shot autoformalization tasks.

\item \textbf{Natural Language Understanding, Reasoning, and Code}: To thoroughly evaluate general linguistic, reasoning, and programming proficiencies, DeepSeekMath-Base is benchmarked on the Massive Multitask Language Understanding (MMLU) dataset \citep{mmlu}, which features 57 multiple-choice tasks across a range of disciplines; BIG-Bench Hard (BBH) \citep{bbh}, a suite of 23 complex problems primarily necessitating multi-step reasoning; and the HumanEval \citep{codex} and MBPP \citep{mbpp} datasets, both standard for code model assessment. The incorporation of mathematical pre-training improves both the model’s linguistic comprehension and reasoning abilities.
\end{itemize}

\section{Math Pre-Training}

\subsection{Data Collection and Decontamination} 

This section describes the methodology employed for creating the DeepSeekMath Corpus using Common Crawl resources. As illustrated in Figure \ref{fig:data_collect}, we introduce a stepwise process for systematically extracting a large-scale math dataset from Common Crawl, commencing from a seed set comprising high-quality mathematical data. This pipeline can be similarly adapted for other specialized domains, such as programming.

Initially, the OpenWebMath collection \citep{openwebmath}, recognized for its well-curated mathematical web content, is selected as the seed dataset. Leveraging this collection, a fastText classifier \citep{joulin2016fasttext} is trained to retrieve additional web pages exhibiting characteristics similar to OpenWebMath. For model training, 500,000 samples are randomly taken from the seed as positive instances, and another 500,000 web pages from Common Crawl as negative instances. The training procedure utilizes an open-source implementation\footnote{\url{https://fasttext.cc}}, with parameters set as follows: 256-dimensional vectors, a learning rate of 0.1, a word n-gram length limit of 3, a word frequency minimum of 3, and 3 epochs. To decrease the dataset size from Common Crawl, both URL-based deduplication and near-duplicate filtering are performed, narrowing down the corpus to 40 billion HTML web pages. Utilizing the trained fastText classifier, math-relevant web pages are retrieved from this deduplicated set. To exclude material of insufficient quality, web pages are sorted based on the scores assigned by the fastText model, retaining only the highest-ranked examples. The retained dataset size is determined through pre-training experiments involving subsets of 40B, 80B, 120B, and 160B tokens. For the first cycle, only the top 40B tokens are preserved.

Following the initial round of data collection, a significant portion of mathematical web content remained uncaptured, primarily due to the insufficient diversity in the positive samples used to train the fastText classifier. To address this, we broadened the seed dataset by incorporating additional mathematical web domains, thereby enhancing the performance of the fastText model. Concretely, the Common Crawl was partitioned into non-overlapping domains, each defined by a unique base URL. For each domain, we determined the fraction of web pages collected during the first pass; any domain with more than 10\% of its pages collected was labeled as math-focused (for instance, \url{mathoverflow.net}). 

Next, we manually reviewed and annotated URLs within these domains that were indicative of mathematical content (such as \url{mathoverflow.net/questions}). Uncollected web pages corresponding to these URLs were then included in the seed set. This iterative expansion of the positive example set allowed for improved training of the fastText model, which in turn facilitated broader retrieval of mathematical content in subsequent iterations. After completing four rounds of collection, we amassed a total of 35.5M mathematical web pages, equivalent to 120B tokens. By the fourth iteration, we observed that almost 98\% of the data had already been obtained by the previous round, leading us to halt further collection.

To mitigate the risk of benchmark data contamination, we adopted the filtering procedures outlined by \cite{deepseek-coder}. Specifically, we excluded from our math training dataset any web pages containing question or answer content from English mathematical benchmarks like GSM8K \citep{gsm8k} and MATH \citep{MATH}, as well as Chinese benchmarks such as CMATH \citep{wei2023cmath} and AGIEval \citep{agieval}. Our filtering method discards any text segment where a 10-gram sequence precisely matches any substring from the benchmark datasets. For benchmark passages that are fewer than 10 but at least 3 tokens long, we removed web pages via exact string match.

\subsection{Assessing the Quality of the DeepSeekMath~Corpus}

To evaluate the quality of the DeepSeekMath~Corpus, we conducted pre-training experiments, comparing it to several recent large-scale mathematical text corpora:

\begin{itemize}[topsep=0pt]
    \item \textbf{MathPile} \citep{mathpile}: a composite dataset (8.9B tokens) aggregated from various sources including textbooks, Wikipedia, ProofWiki, CommonCrawl, StackExchange, and arXiv, with over 85\% originating from arXiv;
    \item \textbf{OpenWebMath} \citep{openwebmath}: a subset of CommonCrawl filtered to retain mathematical content, consisting of 13.6B tokens;
    \item \textbf{Proof-Pile-2} \citep{llemma}: a mathematical compilation comprising OpenWebMath, AlgebraicStack (10.3B tokens of mathematical code), and arXiv papers (28.0B tokens). In our experiments using Proof-Pile-2, we adopted the arXiv:Web:Code ratio of 2:4:1, as recommended in \cite{llemma}.
\end{itemize}

\subsubsection{Training Setting}
\label{sec:quality-policy}
We fine-tune a general pre-trained language model consisting of 1.3B parameters—architecturally equivalent to the DeepSeek LLMs \citep{deepseek-llm}, and hereafter referred to as DeepSeek-LLM 1.3B—specifically for mathematics-related tasks. Each distinct mathematical dataset is used to individually train the model over 150B tokens. Training procedures utilize the lightweight and efficient HAI-LLM \citep{haillm} framework. Consistent with the optimization protocol of DeepSeek LLMs, we adopt the AdamW optimizer \citep{adamW}, setting $\beta_1=0.9$, $\beta_2=0.95$, and $\mathrm{weight\_decay}=0.1$. The learning rate follows a multi-stage schedule: it peaks following 2,000 warmup iterations, is reduced to 31.6\% of the peak at the 80\% mark of training, and further declines to 10.0\% at the 90\% mark. The peak learning rate is capped at $5.3\times 10^{-4}$. Training batches consist of 4 million tokens, with each context window spanning 4,000 tokens.

\subsubsection{Evaluation Results}

\textbf{The DeepSeekMath~Corpus demonstrates superior quality, substantial multilingual coverage, and exceeds existing datasets in scale.}

\begin{itemize}[topsep=0pt]
    \item \textbf{High-quality}:
    To assess downstream capabilities, we conduct evaluations across 8 mathematical benchmarks utilizing few-shot chain-of-thought prompting \cite{cot}. Table \ref{tab:corpora_comparison} illustrates that models trained on the DeepSeekMath~Corpus consistently achieve the highest scores. Additionally, Figure \ref{fig:corpora_comparisons} reveals that, even at 50B tokens (corresponding to a complete epoch of Proof-Pile-2), DeepSeekMath-trained models outperform those trained on Proof-Pile-2, suggesting that DeepSeekMath offers higher overall data quality. 
    \item \textbf{Multilingual}:
    The DeepSeekMath~Corpus incorporates material from several languages, with English and Chinese forming the predominant share. Table \ref{tab:corpora_comparison} indicates that using DeepSeekMath~Corpus for training leads to notable advances in mathematical reasoning for both English and Chinese, whereas other available datasets—which are primarily English-focused—fail to enhance and sometimes even degrade Chinese language performance.
    \item \textbf{Large-scale}:
    The DeepSeekMath~Corpus dwarfs prior mathematical corpora in terms of sheer data volume. As indicated in Figure \ref{fig:corpora_comparisons}, models trained on DeepSeekMath~Corpus not only improve more rapidly, but also sustain these gains over a longer period. Conversely, existing baseline corpora, limited in size, require multiple passes to exhaust, resulting in model performance that quickly saturates.
\end{itemize}

\subsection{Training and Evaluating DeepSeekMath-Base 7B}

This section details the development of DeepSeekMath-Base 7B, a foundational model characterized by robust mathematical reasoning capabilities. The model is initialized using DeepSeek-Coder-Base-v1.5 7B \citep{deepseek-coder} and subsequently trained on a total of 500B tokens. The training dataset is composed of 56\% DeepSeekMath Corpus, 4\% AlgebraicStack, 10\% arXiv materials, 20\% Github code, and 10\% natural language texts sourced from Common Crawl, encompassing both English and Chinese. The majority of training configurations follow those described in Section \ref{sec:quality-policy}, with two key adjustments: the peak learning rate is increased to 4.2e-4 and the batch size is set to 10M tokens.

A thorough evaluation is conducted to analyze the mathematical proficiency of DeepSeekMath-Base 7B, emphasizing the model’s ability to generate self-contained mathematical solutions without external assistance, utilize tools for mathematical problem-solving, and perform formal theorem proving. In addition, the model’s capabilities are profiled more broadly, including assessments of its skills in natural language understanding, logical reasoning, and programming.

\paragraph{Mathematical Problem Solving with Step-by-Step Reasoning}

DeepSeekMath-Base's aptitude for mathematical problem-solving is tested using few-shot chain-of-thought prompting \citep{cot} across eight different benchmarks in both English and Chinese. These benchmarks evaluate abilities in quantitative reasoning (for example, GSM8K \citep{gsm8k}, MATH \citep{MATH}, and CMATH \citep{wei2023cmath}) and multiple-choice problem sets (such as MMLU-STEM \citep{mmlu} and Gaokao-MathQA \citep{agieval}), representing a range of mathematical domains from elementary topics through college-level material.

As detailed in Table \ref{tab:base_math_cot}, DeepSeekMath-Base 7B consistently achieves top results across all eight evaluated benchmarks when compared to other open-source base models, such as the commonly adopted Mistral 7B \citep{mistral} and the newly introduced Llemma 34B \citep{llemma}, which was further trained on Proof-Pile-2 \citep{llemma} for mathematical expertise. Particularly noteworthy is DeepSeekMath-Base’s performance on the MATH benchmark—focused on competition-level problems—where it exceeds the best open-source base models by more than 10 percentage points and even surpasses Minerva 540B \citep{minerva}. Minerva, a closed-source base model 77 times larger and developed atop PaLM \citep{palm} with additional math-centric training, lags behind DeepSeekMath-Base in this regard.

\paragraph{Mathematical Problem Solving with Tool Use} 
For evaluating the models’ ability to perform program-assisted mathematical reasoning, we apply few-shot program-of-thought prompting \citep{pot,pal} to GSM8K and MATH. The models are instructed to address each question by constructing a Python script, making use of libraries such as \textit{math} and \textit{sympy} for advanced calculations. The output of each script serves as the predicted answer. Table \ref{tab:base_math_pot_proof} illustrates that DeepSeekMath-Base 7B achieves state-of-the-art results, outperforming the previous leading Llemma 34B.

\paragraph{Formal Mathematics}

Automating the formalization of proofs plays a pivotal role in enhancing both the accuracy and efficiency of mathematical reasoning and has received growing attention recently. DeepSeekMath-Base 7B is assessed on the informal-to-formal proof task introduced in \citep{dsp_proof}, which involves producing a formal proof based on an informal theorem statement, its formal version, and a corresponding informal proof. We conduct experiments on the miniF2F benchmark \citep{minif2f}, which encompasses formalized Olympiad-level mathematics, and use few-shot prompting to generate formal proofs in Isabelle for each problem. Adopting the method from \cite{dsp_proof}, our models are tasked with creating proof sketches, after which we apply the Sledgehammer automated prover \citep{sledgehammer} to complete any remaining proof steps. As demonstrated in Table \ref{tab:base_math_pot_proof}, DeepSeekMath-Base 7B displays notable performance in automated proof formalization.

\paragraph{Natural Language Understanding, Reasoning, and Code}

Model competence in natural language understanding is measured using MMLU \citep{mmlu}, reasoning through BBH \citep{bbh}, and programming via HumanEval \citep{codex} and MBPP \citep{mbpp}. Table \ref{tab:understanding_reasoning_code} reveals that DeepSeekMath-Base 7B provides a clear improvement on MMLU and BBH compared to its predecessor, DeepSeek-Coder-Base-v1.5 \citep{deepseek-coder}, highlighting the benefits of specialized mathematical training for language and reasoning tasks. Additionally, incorporating code-related tokens into ongoing training allows DeepSeekMath-Base 7B to preserve the coding proficiency of DeepSeek-Coder-Base-v1.5 across HumanEval and MBPP. Overall, DeepSeekMath-Base 7B achieves substantially better results than Mistral 7B \citep{mistral} on all three reasoning and programming evaluations.

\section{Supervised Fine-Tuning}

\subsection{SFT Data Curation}

We compile a mathematical instruction-tuning dataset encompassing both English and Chinese questions from a range of mathematical disciplines and at various levels of difficulty. Each problem is paired with solutions presented in chain-of-thought (CoT) \citep{cot}, program-of-thought (PoT) \citep{pot,pal}, or tool-integrated reasoning format \citep{tora}. This curation results in 776K training samples.

\begin{itemize}[topsep=0pt]
    \item \textbf{English mathematical datasets}: Tool-integrated solutions are manually provided for GSM8K and MATH datasets, and we include a selection from MathInstruct \citep{MathInstruct}, in addition to the training partition of Lila-OOD \citep{lila}, featuring problems resolved with either CoT or PoT approaches. This English dataset suite represents a wide array of topics, such as algebra, probability, number theory, calculus, and geometry.
    \item \textbf{Chinese mathematical datasets}: We assemble K-12 mathematics questions in Chinese across 76 sub-topics—for instance, linear equations—and annotate them with both CoT and tool-integrated formats.
\end{itemize}

\subsection{Training and Evaluating DeepSeekMath-Instruct 7B}

We present DeepSeekMath-Instruct 7B, which is further refined through mathematical instruction-tuning based on DeepSeekMath-Base. During training, individual samples are concatenated randomly to form sequences up to a 4K token context window. The training regime consists of 500 steps using a batch size of 256, maintaining a constant learning rate of 5e-5 throughout.

To assess the model’s quantitative reasoning, both with and without tool usage, we utilize four benchmarks covering mathematical tasks in English and Chinese. Comparisons are drawn with leading contemporary models, including:

\begin{itemize}[topsep=0pt]
    \item \textbf{Closed-source models} encompass: (1) members of the GPT lineup, most notably GPT-4 \citep{gpt4} and GPT-4 Code Interpreter~\footnote{\url{https://openai.com/blog/chatgpt-plugins\#\#code-interpreter}}, (2) Gemini Ultra and Pro \citep{gemini}, (3) Inflection-2 \citep{inflection-2}, (4) Grok-1~\footnote{\url{https://x.ai/model-card}}, as well as models from Chinese organizations such as (5) Baichuan-3~\footnote{\url{https://www.baichuan-ai.com}}, and (6) the most recent version of GLM-4~\footnote{\url{https://open.bigmodel.cn/dev/api\#glm-4}} from the GLM family \citep{glm}.
\end{itemize}

The models discussed here are designed for broad applicability and have typically been subjected to various alignment methodologies.
\item \textbf{Open-source models} comprise: general-purpose models such as (1) DeepSeek-LLM-Chat 67B \citep{deepseek-llm}, (2) Qwen 72B \citep{qwen}, (3) SeaLLM-v2 7B \citep{seallm}, and (4) ChatGLM3 6B \citep{chatglm3}; as well as models targeting mathematical proficiency, including (5) InternLM2-Math 20B~\footnote{\url{https://github.com/InternLM/InternLM-Math}}, which builds upon InternLM2 and is subject to dedicated mathematical training and subsequent instruction tuning, (6) Math-Shepherd-Mistral 7B, which leverages PPO optimization \citep{schulman2017proximal} applied to Mistral 7B \citep{mistral} in conjunction with a process-supervised reward mechanism, (7) the WizardMath suite \citep{wizardmath}, which boosts mathematical reasoning capabilities in Mistral 7B and Llama-2 70B \citep{llama2} by adopting evolve-instruct (a variant of instruction tuning utilizing AI-generated instructions) and PPO training with primary data from GSM8K and MATH, (8) MetaMath 70B \citep{metamath}, which is Llama-2 70B fine-tuned on an expanded GSM8K and MATH corpus, (9) ToRA 34B \cite{tora}, a CodeLlama 34B model refined for tool-augmented mathematical reasoning, and (10) MAmmoTH 70B \citep{MathInstruct}, which involves instruction tuning of Llama-2 70B on MathInstruct.
\end{itemize}

According to Table \ref{tab:sft_rl_math}, in the evaluation context prohibiting tool assistance, DeepSeekMath-Instruct 7B demonstrates exceptional capability in stepwise reasoning. Remarkably, on the advanced MATH dataset, our model outperforms every open-source competitor and nearly all closed-source systems (such as Inflection-2 and Gemini Pro) by no less than 9 percentage points in absolute terms. This holds even when compared to much larger models (such as Qwen 72B) or those specifically reinforced for mathematics (e.g., WizardMath-v1.1 7B). While DeepSeekMath-Instruct’s results are on par with Chinese proprietary models like GLM-4 and Baichuan-3 on the MATH benchmark, it still trails GPT-4 and Gemini Ultra.

When shifting to an evaluation setup that permits the use of both natural language reasoning and programmatic tools, DeepSeekMath-Instruct 7B achieves accuracy close to 60\% on the MATH dataset, exceeding all currently available open-source alternatives. On other benchmarks, our system is on par with DeepSeek-LLM-Chat 67B, a leading model ten times its scale.
\section{Reinforcement Learning}

\subsection{Group Relative Policy Optimization} Reinforcement learning (RL) has demonstrated its value for further enhancing the mathematical reasoning skills of large language models after the Supervised Fine-Tuning (SFT) phase \citep{wang2023math,wizardmath}. In this part, we present Group Relative Policy Optimization (GRPO), our streamlined and potent RL approach.

\subsubsection{From PPO to GRPO} Proximal Policy Optimization (PPO) \citep{schulman2017proximal}, an actor-critic RL algorithm, is widely adopted in the RL fine-tuning process for LLMs \citep{ouyang2022training}. The algorithm works by maximizing the following surrogate objective to optimize LLM parameters:

\begin{equation}
\footnotesize
    \mathcal{J}_{PPO}(\theta) = \mathbb{E}{[q \sim P(Q), o \sim \pi_{\theta_{old}}(O|q)]} \frac{1}{|o|} \sum_{t=1}^{|o|} \min \left[ \frac{\pi_\theta(o_{t} | q, o_{<t})}{\pi_{\theta_{old}}(o_{t} | q, o_{<t})} A_{t}, \text{clip} \left( \frac{\pi_\theta(o_{t} | q, o_{<t})}{\pi_{\theta_{old}}(o_{t} | q, o_{<t})}, 1 - \varepsilon, 1 + \varepsilon \

Here, $\pi_{\theta}$ and $\pi_{\theta_{old}}$ denote the current and preceding policy networks, while $q$ and $o$ correspond to questions and their respective outputs, sampled both from the question dataset and from the historical policy $\pi_{\theta_{old}}$. The parameter $\varepsilon$ is a PPO-specific clipping hyperparameter designed to promote stable training dynamics. The advantage term $A_t$ is derived using Generalized Advantage Estimation (GAE) \citep{gae}, leveraging the future rewards $\{r_{\ge t}\}$ and a parameterized value function $V_{\psi}$. Consequently, PPO requires training a value network in parallel with the policy network. To counteract excessive optimization of the reward model, it is standard to append a per-token KL divergence penalty—measured against a reference model—to each reward value \citep{ouyang2022training}, as follows:

\begin{equation}
    r_{t} = r_\varphi(q, o_{\le t}) - \beta \log\frac{\pi_{\theta}(o_{t}|q, o_{<t})}{\pi_{ref}(o_{t}|q, o_{<t})},
\label{eq:PPO-reward}
\end{equation}

In this equation, $r_\varphi$ is the learned reward model, $\pi_{ref}$ is the reference distribution (commonly the initial SFT model), and $\beta$ tunes the weight of the KL penalty.

However, because PPO’s value function often has model complexity on par with the policy network, this leads to significant memory and computation costs. Furthermore, during RL optimization, the value function serves as a baseline in the calculation of the advantage in order to suppress variance. In the setting of LLMs, the reward model is typically applied only to the final token of a sequence, complicating the process of learning accurate tokenwise value predictions. To mitigate these issues, as illustrated in Figure \ref{fig:grpo}, we introduce Group Relative Policy Optimization (GRPO), which removes the necessity for a separate value function, as in PPO. Instead, GRPO relies on the mean reward obtained from multiple sampled completions generated for the same question, using this average as a baseline. Concretely, for every question $q$, GRPO samples a set of outputs $\{o_1, o_2, \cdots, o_G\}$ from $\pi_{\theta_{old}}$ and subsequently updates the policy by maximizing the following criterion:

\begin{equation}
\footnotesize
\begin{split}
    \mathcal{J}_{GRPO}(\theta) &= \mathbb{E}{[q \sim P(Q), \{o_i\}_{i=1}^G \sim \pi_{\theta_{old}}(O|q)]}  \\
    & \frac{1}{G}\sum_{i=1}^G\frac{1}{|o_i|} \sum_{t=1}^{|o_i|} \left\{ \min \left[ \frac{\pi_\theta(o_{i,t} | q, o_{i,<t})}{\pi_{\theta_{old}}(o_{i,t} | q, o_{i,<t})} \hat{A}_{i,t}, \text{clip} \left( \frac{\pi_\theta(o_{i,t} | q, o_{i,<t})}{\pi_{\theta_{old}}(o_{i,t} | q, o_{i,<t})}, 1 - \varepsilon, 1 + \varepsilon \right)  \hat{A}_{i,t} \right] - \beta \mathbb{D}_{KL}\left[\pi_{\theta} || \pi_{ref}\right]\right\} ,
\end{split}
\label{eq:GRPO-obj}
\end{equation}

Here, $\varepsilon$ and $\beta$ continue to serve as hyperparameters, and the advantage $\hat{A}_{i,t}$ is defined exclusively with respect to the relative rewards among the outputs in the sampled group—specifics of which are elaborated in subsequent sections. By calculating advantages comparatively within groups, GRPO capitalizes on the fact that reward models are often built on the basis of pairwise comparisons for responses to the same prompt. Furthermore, in contrast to inserting the KL penalty directly in the reward term as in PPO, GRPO regularizes by adding the KL divergence between the policy under training and the reference model directly to the objective function, thus simplifying the computation of $\hat{A}_{i,t}$. Distinct from the KL penalty in (\ref{eq:PPO-reward}), the KL divergence is estimated with the following unbiased estimator \

\begin{algorithm}[t]
  \small
  \caption{Iterative Group Relative Policy Optimization}
  \textbf{Input} initial policy model $\pi_{\theta_{\text{init}}}$; reward models $r_\varphi$; task prompts $\mathcal{D}$; 
  hyperparameters $\varepsilon$, $\beta$, $\mu$
  \begin{algorithmic}[1]
    \State Initialize policy model: $\pi_\theta \leftarrow \pi_{\theta_{\text{init}}}$
    \For{iteration = 1, \dots, I}
      \State Set reference model $\pi_{ref} \leftarrow \pi_{\theta}$
      \For{step = 1, \dots, M}
        \State Randomly sample a batch $\mathcal{D}_b$ from $\mathcal{D}$
        \State Assign previous policy $\pi_{\theta_{old}} \leftarrow \pi_{\theta}$
        \State For each question $q \in \mathcal{D}_b$, generate $G$ outputs $\{o_i\}_{i=1}^G$ using $\pi_{\theta_{old}} (\cdot \mid q)$
        \State Use $r_{\varphi}$ to evaluate each $o_i$ and obtain corresponding rewards $\{r_i\}_{i=1}^G$
        \State Estimate group relative advantage $\hat{A}_{i,t}$ for each token $t$ in output $o_i$
        \For{GRPO iteration = 1, \dots, $\mu$}
          \State Update $\pi_{\theta}$ via maximization of the GRPO objective (Equation \ref{eq:GC-GRPO})
        \EndFor
      \EndFor
      \State Adapt $r_\varphi$ through ongoing training and experience replay. 
    \EndFor 
  \end{algorithmic}
  \textbf{Output} $\pi_\theta$
  \label{alg:iter-grpo}
\end{algorithm}

\subsubsection{Outcome Supervision RL with GRPO} In a formal sense, for any question $q$, one draws a set of $G$ responses $\{o_1, o_2, \cdots, o_G\}$ from the preceding policy $\pi_{\theta_{old}}$. These responses are evaluated using a reward model, producing a vector of rewards $\mathbf{r} = \{r_1, r_2, \cdots, r_G\}$. The rewards undergo normalization by deducting the group mean and scaling by the group’s standard deviation. Under outcome supervision, the resulting normalized score is delivered only at the completion of each response $o_i$, and this value is propagated as the advantage $\hat{A}_{i, t}$ for all tokens within $o_i$: $\hat{A}_{i, t} = \widetilde{r}_i = \frac{r_i- {\rm mean}(\mathbf{r})}{{\rm std}(\mathbf{r})}$. Policy optimization is then carried out by maximizing the objective presented in equation (\ref{eq:GRPO-obj}).

\subsubsection{Process Supervision RL with GRPO} Relying solely on terminal rewards—as in outcome supervision—often falls short when guiding policies for sophisticated mathematical reasoning. Drawing on insights from \cite{wang2023math}, we investigate process supervision, in which each reasoning step is individually scored. Specifically, for question $q$ with $G$ sampled outputs $\{o_1, o_2, \cdots, o_G\}$, a process-level reward model assigns rewards for each step, resulting in: $\mathbf{R} = \{ \{r_1^{index(1)},\cdots,r_1^{index(K_1)}\}, \cdots,  \{r_G^{index(1)},\cdots,r_G^{index(K_G)}\} \}$. Here, $index(j)$ corresponds to the end token position of step $j$, and $K_i$ denotes the total reasoning steps in output $o_i$. These process-level rewards are normalized by subtracting the aggregate mean and dividing by the standard deviation: $\widetilde{r}_i^{index(j)} = \frac{r_i^{index(j)} - {\rm mean}(\mathbf{R})}{{\rm std}(\mathbf{R})}}$. Afterwards, the advantage for token $t$ is computed as the cumulative normalized rewards across all subsequent reasoning steps, i.e., $\hat{A}_{i, t} = \sum_{index(j) \ge t} \widetilde{r}_i^{index(j)}$, and the policy is refined through maximization

\subsubsection{Iterative RL with GRPO} During the reinforcement learning process, earlier versions of the reward model may no longer provide adequate supervision for the evolving policy model. To address this, we also investigate an iterative RL framework using GRPO. As depicted in Algorithm \ref{alg:iter-grpo}, iterative GRPO involves producing updated datasets for training the reward model by sampling outputs from the policy model. The reward model is continually refined using a replay strategy that maintains 10\% of previously collected data. Subsequently, the policy model is set as the new reference model, and its training proceeds under the guidance of the updated reward model.

\subsection{Training and Evaluating DeepSeekMath-RL}

Our reinforcement learning experiments employ DeepSeekMath-Instruct 7B as the foundation. The RL training set is composed of chain-of-thought questions drawn from GSM8K and MATH within the SFT dataset, amounting to approximately 144K questions. In order to analyze the effect of RL on benchmarks not represented in the RL training phase, we exclude all other SFT examples. The construction of the reward model’s training set follows the methodology presented by \citep{wang2023math}. The initial reward model is trained starting from DeepSeekMath-Base 7B, using a learning rate of 2e-5. For the GRPO method, we set the learning rate of the policy model at 1e-6 and adopt a KL coefficient of 0.04. We generate $64$ candidate completions per question, with a maximum sequence length of 1024 tokens and a training batch size of 1024. The policy model receives one update after each round of exploration. DeepSeekMath-RL 7B is evaluated using the same suite of benchmarks as DeepSeekMath-Instruct 7B. For DeepSeekMath-RL 7B, the GSM8K and MATH datasets featuring chain-of-thought reasoning are treated as in-domain tasks, whereas all other benchmark datasets are considered out-of-domain.

Table \ref{tab:sft_rl_math} provides a comparative analysis of both open- and closed-source models, evaluating their performance on English and Chinese benchmarks utilizing chain-of-thought and tool-augmented reasoning methods. Our observations are as follows: 1) When employing chain-of-thought reasoning, DeepSeekMath-RL 7B achieves 88.2\% accuracy on GSM8K and 51.7\% on MATH, outperforming every open-source model in the 7B to 70B parameter category and most closed-source alternatives. 2) Notably, the training of DeepSeekMath-RL 7B relies exclusively on chain-of-thought-style instruction tuning data for GSM8K and MATH, and is initialized from DeepSeekMath-Instruct 7B. Despite its narrowly focused training set, DeepSeekMath-RL 7B exceeds the performance of DeepSeekMath-Instruct 7B across all assessed metrics, underlining the value of reinforcement learning.

\section{Discussion}
Here, we present insights derived from our pre-training and reinforcement learning experiments.

\subsection{Lessons Learnt in Pre-Training}

We begin by sharing our observations on pre-training. Unless explicitly stated, the training protocols align with those described in Section \ref{sec:quality-policy}. For the purposes of this discussion, references to the DeepSeekMath Corpus pertain to an 89B-token dataset compiled during the second cycle of data collection.

\subsubsection{Code Training Benefits Mathematical Reasoning}

An often-cited but insufficiently substantiated claim is that code-centric training enhances reasoning capabilities. Our experiments provide evidence—specifically within the mathematical context—that training on code boosts a model's proficiency in mathematical reasoning, both in scenarios with and without auxiliary tool usage.

To evaluate the influence of code training on mathematical reasoning skills, we investigated the following training regimes: two-stage training and one-stage training configurations.

\textbf{Two-Stage Training}

\begin{itemize}[topsep=0pt]
    \item \textbf{Code Training for 400B Tokens $\rightarrow$ Math Training for 150B Tokens}: DeepSeek-LLM 1.3B undergoes pre-training on 400B code tokens, followed by an additional 150B tokens consisting of math data.
    \item \textbf{General Training for 400B Tokens $\rightarrow$ Math Training for 150B Tokens}: As a comparison baseline, we perform the initial training phase with general language tokens—sampled from DeepSeek-AI’s extensive general corpus—instead of code. This design enables examination of whether pre-training on code, relative to generic language data, confers unique benefits for mathematical reasoning skills.
\end{itemize}

\textbf{One-Stage Training}

\begin{itemize}[topsep=0pt]
    \item \textbf{Math Training for 150B Tokens}: DeepSeek-LLM 1.3B is trained directly on a 150B-token math dataset.
    \item \textbf{Training on a mixture of 400B Code Tokens and 150B Math Tokens}: Since continued math training after code training tends to impair coding abilities, we probe whether jointly training on a composite set of code and math tokens in a single stage can simultaneously enhance mathematical reasoning and lessen the detrimental effects of catastrophic forgetting.
\end{itemize}

\paragraph{Results}
Table \ref{tab:code-math} and Table \ref{tab:code-math-general-eval} summarize downstream outcomes for these distinct training regimes.

Integrating code training fosters improvements in program-augmented mathematical reasoning, observable in both two-stage and single-stage paradigms. Specifically, as illustrated in Table \ref{tab:code-math}, a code-only pre-training phase within the two-stage framework already significantly elevates the model’s proficiency at solving GSM8K and MATH datasets via Python solutions. An additional math-specialized fine-tuning phase further enhances these outcomes. Notably, when code and math tokens are combined in a single-stage process, this approach not only effectively counters the problem of catastrophic forgetting inherent to sequential stagewise training, but also simultaneously promotes both code-based performance (see Table \ref{tab:code-math-general-eval}) and program-supported mathematical reasoning (see Table \ref{tab:code-math}).

Pre-training on code also bolsters mathematical problem-solving absent tool augmentation. With the two-stage schedule, initial exposure to code yields moderate performance gains on reasoning tasks, while subsequent math training is rendered more effective, together achieving superior results. In contrast, a single-stage mixture of code and math tokens appears to hinder unaided mathematical reasoning. A plausible explanation for this phenomenon is that the relatively modest scale of DeepSeek-LLM 1.3B constrains its representational capacity, precluding full mastery of both code and math domains within a joint training configuration.

\subsubsection{ArXiv Papers Appear to Contribute Little to Advancing Mathematical Reasoning}

Although arXiv papers are frequently used in constructing math-focused pre-training corpora \citep{minerva,gpt-f,llemma,mathpile}, there is a lack of in-depth investigation into their actual influence on mathematical reasoning performance. Somewhat unexpectedly, our experimental evidence suggests that arXiv papers do not significantly enhance mathematical reasoning abilities. We carried out systematic experiments utilizing various model sizes, such as DeepSeek-LLM 1.3B and DeepSeek-Coder-Base-v1.5 7B \citep{deepseek-coder}, training them on arXiv-based datasets processed by different methods:

\begin{itemize}[topsep=0pt]
    \item \textbf{MathPile} \citep{mathpile}: This corpus, sized at 8.9B tokens, is curated using cleaning and heuristic filtering, with over 85\% comprised of scientific articles from arXiv;
    \item \textbf{ArXiv-RedPajama} \citep{redpajama}: Consists of all arXiv LaTeX source files after eliminating preambles, commentary, macro definitions, and reference lists, amassing 28.0B tokens.
\end{itemize}

Our experiments involved independently training DeepSeek-LLM 1.3B for 150B tokens and DeepSeek-Coder-Base-v1.5 7B for 40B tokens on each respective arXiv dataset. Results consistently indicate that exposure to arXiv content alone offers negligible—if any—gains in mathematical reasoning, and in certain cases, leads to a degradation in performance, as measured by a diverse set of mathematical benchmarks of varying difficulty used in this analysis. The evaluation suite includes numerical problem solving datasets such as GSM8K and MATH (Table \ref{tab:arxiv-cot}), standardized multiple-choice tasks like MMLU-STEM (Table \ref{tab:arxiv-cot}), and formal mathematical proof datasets such as miniF2F (Table \ref{tab:arxiv_atp}).

It should be emphasized, nevertheless, that this finding is bounded by several unaddressed aspects. Our current study does not explore:

\begin{itemize}[topsep=0pt]
    \item The possible effects of arXiv-derived tokens on specific mathematical subtasks excluded from our benchmarks, for example, the informalization process of translating formalized statements or proofs into natural language;
    \item The role of arXiv data when it is combined with other sources within a multi-corpus training setting;
    \item Potential advantages that might arise at larger model scales beyond those tested.
\end{itemize}

Hence, these open questions merit further investigation, which we defer to subsequent research.

\subsection{Insights of Reinforcement Learning}
\subsubsection{Toward a Unified Framework} This section introduces a general framework to interpret and compare several training strategies, namely SFT, RFT, DPO, PPO, and GRPO. We also carry out experiments probing the underlying elements of this unified framework. The gradient of a training objective with respect to the parameter $\theta$ for these methods can typically be expressed as:

\begin{equation}
    \nabla_{\theta}\mathcal{J}_{\textcolor{red}{\mathcal{A}}}(\theta) = \mathbb{E}[\underbrace{(q,o) \sim \textcolor{red}{\mathcal{D}}}_{Data \ Source}]\left( \frac{1}{|o|} \sum_{t=1}^{|o|}  \underbrace{GC_{{\mathcal{A}}}(q, o, t, \textcolor{red}{\pi_{{rf}}})}_{Gradient \ Coefficient}  \nabla_{\theta}\log \pi_{\theta}(o_t | q, o_{<t})\right).
\label{eq:objective}
\end{equation}

This framework involves three main components: (1) \textit{Data Source $\mathcal{D}$}, defining the set of data used for training; (2) \textit{Reward Function $\pi_{{rf}}$}, supplying the evaluative feedback necessary for guiding the model; (3) \textit{Algorithm $\mathcal{A}$}, which incorporates both the training samples and reward signals to compute the gradient coefficient $GC$, quantifying the extent of reward or penalty imparted to specific samples. We explore several key approaches formulated within this shared structure:

\begin{itemize}
    \item  \textbf{Supervised Fine-tuning (SFT)}: This strategy adapts a pretrained model through additional training on SFT datasets curated by humans.
    \item \textbf{Rejection Sampling Fine-tuning (RFT)}: Building upon the SFT model, RFT continues training using outputs that have been filtered for answer correctness after being generated in response to SFT prompts.
    \item \textbf{Direct Preference Optimization (DPO)}: DPO extends the SFT approach by optimizing on enriched outputs derived from the SFT model using a pairwise DPO loss to drive learning.
    \item \textbf{Online Rejection Sampling Fine-tuning (Online RFT)}: Unlike RFT, Online RFT updates the policy model (initialized from the SFT model) by continuously fine-tuning it with outputs generated from the current policy itself.
    \item \textbf{PPO/GRPO}: With PPO/GRPO, the SFT model is used as a starting point, after which the model is further trained and rewarded with outputs sampled dynamically from the current version of the policy.
\end{itemize}

We present a summary of these method components in Table \ref{tab:data_policy}. For a thorough derivation, please consult Appendix \ref{app:analysis-rl}.

\paragraph{Observation about Data Source} We categorize data sources into online and offline sampling. Online sampling involves generating training data from the ongoing exploration behavior of the current policy model during training. Conversely, offline sampling utilizes training data produced by the initial SFT model before further policy adaptation. Notably, RFT and DPO employ the offline sampling paradigm, whereas Online RFT and GRPO adopt online sampling strategies.

As depicted in Figure \ref{fig:rl-analysis}, our analysis shows that Online RFT delivers a marked performance improvement over RFT on both benchmarks. In the initial phases of training, Online RFT and RFT achieve similar results, but as training progresses, Online RFT establishes a clear lead, underscoring the benefits of online data collection. This trend is expected; early on, both the actor and the SFT model behave similarly, producing samples that differ only slightly. Over time, however, the actor's outputs diverge more substantially from those of the SFT, making the real-time sampled data increasingly valuable.

\paragraph{Observation about Gradient Coefficient} In these algorithms, the reward signal is transformed into a gradient coefficient to direct parameter updates. Our experiments distinguish two approaches to the reward function: 'Rule' and 'Model'. In the 'Rule' setup, response quality is assessed strictly by answer correctness, whereas in the 'Model' approach, a reward model—trained on data labeled according to rule-based judgments—assigns a score to each response. Equations \ref{eq:GC-RFT} and \ref{eq:GC-GRPO} illuminate an essential contrast between GRPO and Online RFT. GRPO introduces a reward-dependent adjustment to the gradient coefficient, modulating reinforcement and penalization in proportion to the reward magnitude assigned by the reward model. In contrast, Online RFT lacks this granularity, treating all correct responses equivalently and refraining from explicitly penalizing incorrect answers.

As illustrated in Figure \ref{fig:rl-analysis}, GRPO outperforms online RFT, thereby underscoring the efficacy of adjusting the coefficients associated with positive and negative gradients. Moreover, GRPO+PS achieves better results than GRPO+OS, demonstrating the advantage conferred by employing step-aware, fine-grained gradient coefficients. We further investigate the impact of iterative RL, running two cycles in our experimental setup. According to Figure \ref{fig:iter-rl}, iterative RL yields noticeable gains, with the first iteration delivering the most pronounced improvement.

\subsubsection{Why RL Works?} This work applies reinforcement learning to a subset of instruction tuning data, leading to substantial improvements over models trained solely with instruction tuning. To shed light on the underlying reasons for RL’s effectiveness, we compare the Pass@K and Maj@K accuracies of Instruct and RL-enhanced models across two benchmark datasets. As depicted in Figure \ref{fig:maj-pass}, RL increases Maj@K scores, yet leaves Pass@K unchanged. These results suggest that RL bolsters general model robustness primarily by promoting correct answers within the TopK set, rather than fundamentally advancing underlying capabilities. Notably, \citep{wang2023making} documented a \textbf{misalignment problem} in reasoning tasks for SFT models, and demonstrated that such reasoning can be improved through various preference alignment methodologies \citep{yuan2023rrhf,song2023preference,wang2023making}.

\subsubsection{How to Achieve More Effective RL?}
\label{sec:effec_rl}
We establish that RL yields substantial benefits for mathematical reasoning tasks. Furthermore, we introduce a unified framework that accommodates diverse, well-known training approaches, framing them as either explicit or simplified RL methods. As formalized in Equation \ref{eq:objective}, three critical factors structure this framework: Data Source, Algorithm, and Reward Function. Below, we propose several potential future avenues concerning each of these components.

\paragraph{Data Source} The data source underpins all training approaches. For RL, this typically means unlabeled questions and the policy model’s corresponding sampled outputs. This work restricts its scope to questions employed during instruction tuning and utilizes basic nucleus sampling to obtain candidate responses. We posit that this may limit improvements exclusively to Maj@K. To address this, future investigations will apply our RL framework to out-of-distribution question prompts and integrate \textbf{advanced sampling (decoding) techniques} such as those employing tree-search \citep{yao2023tree}. Additionally, \textbf{efficient inference methods} \citep{xia-etal-2023-speculative,leviathan2023fast,kwon2023efficient,xia2024unlocking}, which determine the exploration capacity of policy models, are also of critical significance.

\paragraph{Algorithms} Algorithms utilize both the data and the feedback from reward signals to determine the gradient coefficients necessary for updating model parameters. Referring to Equation \ref{eq:objective}, contemporary approaches predominantly \textbf{TRUST} the signals from the reward function to adjust the conditional probability of particular tokens, increasing or decreasing their likelihood accordingly. Nonetheless, these reward signals cannot be guaranteed to be consistently reliable, especially when dealing with highly intricate tasks. For instance, even meticulously annotated resources like the PRM800K dataset \citep{lightman2023let}, prepared by skilled annotators, still include close to 20\% annotation errors\footnote{\url{https://github.com/openai/prm800k/issues/12\#issuecomment-1728491852}}. Given this limitation, our work investigates reinforcement learning algorithms designed to remain effective under noisy or unreliable reward signals. We argue that adopting \textbf{WEAK-TO-STRONG} \citep{burns2023weak} alignment frameworks may lead to transformative shifts in algorithmic learning processes.

\paragraph{Reward Function} The reward function acts as the principal generator of training signals within reinforcement learning frameworks, most commonly instantiated as a neural reward model. We identify three pivotal research directions for such reward models: 1) \textbf{Improving the generalization capacity of reward models.} Reward models should be capable of effectively handling out-of-distribution queries and sophisticated outputs during decoding; otherwise, reinforcement learning risks merely solidifying the existing behavior of large language models (LLMs) instead of genuinely advancing their core capabilities; 2) \textbf{Representing reward model uncertainty.} Accounting for uncertainty can serve as a crucial interface, connecting weak reward models with weak-to-strong learning strategies; 3) \textbf{Developing process-level reward models with efficiency and high fidelity,} such that they can deliver granular training signals to support complex reasoning tasks \citep{lightman2023let,wang2023math}.

\section{Conclusion, Limitation, and Future Work}

In this work, we introduce DeepSeekMath, which establishes a new state-of-the-art among open-source models on the challenging MATH benchmark, and closely rivals the accuracy of proprietary alternatives. DeepSeekMath~originates from DeepSeek-Coder-v1.5 7B and is subsequently trained on 500B tokens, prominently featuring 120B mathematical tokens collected from Common Crawl. Our thorough ablation analyses reveal that web-based resources possess substantial value as high-quality mathematical datasets, whereas the utility of arXiv-sourced data is less pronounced than anticipated. We propose Group Relative Policy Optimization (GRPO), an extension of Proximal Policy Optimization (PPO), which markedly strengthens the model's mathematical reasoning abilities while reducing memory overhead. The empirical evidence indicates that GRPO yields tangible benefits, even when DeepSeekMath-Instruct 7B attains strong performance on existing benchmarks. Moreover, we offer a consolidated framework for interpreting various approaches and identify multiple avenues for advancing reinforcement learning strategies.

Despite DeepSeekMath’s notable achievements in quantitative reasoning, its performance in geometry and proof-based problem-solving remains inferior to that of leading closed-source systems. For example, during preliminary testing, the model demonstrated difficulty with questions involving triangles and ellipses, suggesting that the selection of pre-training and fine-tuning data may introduce biases. Furthermore, due to the model's scale, DeepSeekMath~lags behind GPT-4 in few-shot generalization; unlike GPT-4, which exhibits marked improvements given few-shot examples, DeepSeekMath~shows limited distinction between zero-shot and few-shot settings. Future efforts will focus on refining our data selection methodology to assemble even higher-quality pre-training corpora. Additionally, we intend to further investigate the reinforcement learning approaches outlined in Section \ref{sec:effec_rl} to enhance LLM effectiveness. 

\end{document}